\section{John of the Cross, the Rumi of Spain}

Thus it is perfectly clear how near is \textbf{St. John of the Cross} in his mystical vocabulary (allowing for language differences, of course) and poetic style, i.e., imagery, metaphors, symbolism, even to some extent versification techniques, to the Persian Sufis. One may easily imagine the works of St. John of the Cross without the influence of ibn Abbad, since such influence is found only in \emph{Dark Night of the Soul} and its prose commentary of the same title; however, take away all which so forcefully calls to mind Suhrawardi and the great Persian Sufi poets and St. John of the Cross would be reduced to a spiritual counselor in the tradition of \textbf{Dionysius the Pseudo-Areopagite}, \textbf{St. Gregory of Nyssa} and \textbf{Gregory Palamas} as well as a minor poet, his uniqueness and greatness as poet and thinker gone.

In summary, Persian Sufis and Dervishes came to the Kingdom of Granada, particularly in the 13th and 14th Centuries, and their influence there was very considerable, combining with the Sufi and esoteric traditions of al-Andalus. Later this influence passed to St. John of the Cross by way of the Moriscos and perhaps other means as well. Spain and Safavi Persia were both mortal enemies of the Ottoman Turks; during this time Spaniards did visit Persia, and did learn the Persian language while there. St. John of the Cross may have been in direct or indirect contact with said ambassadors. While there is no documentary proof of this, the thing itself is perfectly possible.

The observations of Fr. Asin relative to the evident influence of ibn Abbad do not weaken one whit the thesis that St. John of the Cross owed a great deal to the Persian Sufis; rather they strengthen it, being yet another proof that St. John of the Cross was in contact with the Moriscos and from them learned much. Nor does all this detract from the genius of St. John of the Cross as a poet and mystic. Of the many millions who have read the verse of the Persian Sufi poets (and St. John of the Cross knew this verse only indirectly), how many have been able to compose such works of genius as are the poems of St. John of the Cross?

The fact that St. John of the Cross was so open to said influence is a proof of the fact that \textbf{he was an initiate in the esoteric path} and of the ecumenism typical of all true mystics. If \textbf{Rumi} is the St. John of the Cross of Persia, then St. John of the Cross is the Rumi of Spain.


\flrightit{Posted on 2016-07-23 by Michael McClain}

\begin{center}* * *\end{center}

\begin{footnotesize}
\begin{sffamily}
\texttt{Allahu Haqq on 2016-07-28 at 22:38 said:}

Where can I get this book?

\hfill

\texttt{Cologero on 2020-07-23 at 07:45 said:}

The text is rambling and needs editing. The author is a technophobe, so the formatting and fonts in the Word document are inconsistent. It needs a lot of work for publication, but the author is uncooperative.

I may just convert everything --- as is --- to PDF files.

\hfill

\texttt{Michael M on 2020-07-23 at 08:25 said:}

If of interest I would be willing to work on the formatting, of course that would take the author being cooperative. PDF as is would be useful anyway for those of us reading.


\hfill

\end{sffamily}
\end{footnotesize}

